\documentclass[10pt,a4paper]{amsart}
\usepackage[margin=19mm]{geometry}
\usepackage{amsmath,amssymb,mathtools}
\usepackage[scr=rsfs,cal=euler]{mathalfa}
\usepackage{xfrac}
\usepackage[unicode,hidelinks,bookmarks=false]{hyperref}
\newcommand{\ie}{i.e.\ }
\newcommand{\cf}{cf.\ }
\newcommand{\resp}{respectively\ }
\newcommand{\Subsection}{\subsection}
\providecommand{\nobreakdash}{\nobreak\hbox{-}}
\setlength{\emergencystretch}{2em}
\multlinegap0pt
\usepackage{bm}
\usepackage{bbm}
\usepackage{dsfont}
\usepackage{xspace}
\usepackage{cancel}
\usepackage{mathtools}
\usepackage{amsthm}
\makeatletter
\@ifundefined{lemma}{\newtheorem{lemma}{Lemma}}{}
\@ifundefined{theorem}{\newtheorem{theorem}{Theorem}}{}
\makeatother
\title[Tilings from Tops of Overlapping Iterated Function Systems]{Tilings from Tops of Overlapping Iterated Function Systems}
\author{Michael F. Barnsley, Corey de Wit}
\date{Source: arXiv:2504.11710v3 (CC BY 4.0)}
\begin{document}
\maketitle

\noindent\emph{Source and licence.} Tilings from Tops of Overlapping Iterated Function Systems. Version arXiv:2504.11710v3 (CC BY 4.0), distributed under the Creative Commons Attribution 4.0 International licence, as recorded in the arXiv licence metadata for this exact version and shown on the abstract page https://arxiv.org/abs/2504.11710v3. Full rights notice: licenses/arxiv-2504.11710-CC-BY-4.0.txt.\par\smallskip

\noindent\textit{Editorial scope.} Counted unit: the Theorem whose source label is \texttt{tiling theorem} in the article (source0.tex lines 416--464, source label \texttt{tiling theorem}), delivered together with the article's own complete original proof of it (source0.tex lines 466--585, one \texttt{proof} environment of 120 lines). No mathematics is written, repaired or re-derived: statement and proof are the article's own text line for line. Required context retained uncounted: coreydef (lines 225--228); topshiftlemma (lines 233--237); smallerlemma (lines 255--259); lemmafour (lines 379--392). Every definition, display and label of those lines is retained unchanged. Omitted from this package and not redistributed: 1--224, 229--232, 238--254, 260--378, 393--415, 465--465, 586--1047. Nothing inside a retained excerpt is omitted or altered: every retained body is delivered exactly as the article's lines are written, so no omission receipt of any kind accompanies this package. Citations: the retained excerpts carry no citation command at all, so no bibliography entry of the article is transcribed and no reference is redistributed. Reference adaptation: none was needed and none was made; every reference inside the delivered unit resolves inside it, so no unresolved reference or citation is produced. Printed slips kept literal: no printed word, symbol, hypothesis or number of the counted statement or of its proof was altered. Layout and class: the article's own class is \texttt{amsproc} with packages including amsfonts, amsmath, amssymb, graphicx, comment; the wrapper is the lane's pinned \texttt{amsart} editorial wrapper at 10pt on A4, which restates the theorem-like environments the retained text needs and adds the article's own macro definitions that the retained text needs (none), with extra wrapper packages (bm, bbm, dsfont, xspace, cancel, mathtools). The delivered numbering is the unit's own. Assets: no retained span contains a figure environment, a graphics-inclusion command, a PDF-page inclusion, a table or a TikZ picture; no image file is copied into this package, the package declares no asset dependency and no image file is redistributed with it. Licence: version arXiv:2504.11710v3 of this article is distributed under the Creative Commons Attribution 4.0 International licence, as recorded in the arXiv licence metadata for this exact version and shown on the abstract page https://arxiv.org/abs/2504.11710v3. A full rights notice is delivered at licenses/arxiv-2504.11710-CC-BY-4.0.txt. Characters: every retained excerpt is pure ASCII, so the delivered engine, XeLaTeX with its default Latin Modern text font, reports no missing character.
\begin{equation}
A_{\mathbf{i}|n}=f_{\mathbf{i}|n}(A)\backslash\cup\{f_{\mathbf{j}%
|n}(A)|\mathbf{j\in}\Sigma,(\mathbf{j}|n)>(\mathbf{i}|n)\}.\label{coreydef}%
\end{equation}

\begin{lemma}
\label{topshiftlemma}Let $\mathbf{k}\in\Sigma.$ Then $f_{k_{1}}\left(
A_{(\sigma\left(  \mathbf{k}\right)  \mathbf{)|}n-1}\right)  \supseteq
A_{\mathbf{k|}n}$ for all $n\in\mathbb{N}$, $n>1$.
\end{lemma}

\begin{lemma}
\label{smallerlemma}Let $n>1.$ If $i_{1}i_{2}\dots i_{n-1}i_{n}\in\Sigma_{n}$,
then both $i_{2}\dots i_{n-1}i_{n}\ $and $i_{1}i_{2}\dots i_{n-1}\ $belong to
$\Sigma_{n-1}.$
\end{lemma}

\begin{lemma}
\label{lemmafour} Let $E$ and $F$ be subsets of $\mathbb{R}^{n}$ such that
$E\backslash F$ is not empty, $E$ is the closure of its interior, and $F$ is
closed. Then the interior $\left(  E\backslash F\right)  ^{\circ}$ of
$E\backslash F$ is not empty.

\begin{proof}
Suppose $\left(  E\backslash F\right)  ^{\circ}$ is empty. Then $E^{\circ
}\subset F.$ Taking the closure of both sides, it follows that $E$ is
contained in $F$. That is, $E\backslash F$ is empty, a contradiction. So if
$E\backslash F$ is not empty, then $\left(  E\backslash F\right)  ^{\circ}$ is
not empty.
\end{proof}
\end{lemma}

\begin{theorem}
\label{tiling theorem} Let $F$ be a strictly contractive IFS on $\mathbb{R}%
^{n}$ whose attractor $A$ has nonempty interior. Let $C$ be the critical set
of $F.$ Let $\mathbf{j\in}\overleftarrow{\Sigma}$, $\widehat{\lambda}%
\in(\lambda,1),$ $x_{0}\in A,$ and $\varepsilon>0,$ be fixed such that the
sequence $\left\{  x_{n}=f_{j_{n}}f_{j_{n-1}}\dots f_{j_{1}}(x_{0})\right\}
_{n=1}^{\infty}$ obeys $d(x_{n},C)>$ $\varepsilon\left(  \widehat{\lambda
}\right)  ^{n}$ for all $n$.

(i) For fixed $l\in\mathbb{N}$ and $\mathbf{i\in}\Omega_{l}$, the sequence
$\left\{  S(l,\mathbf{i},k):k=0,1,\dots\right\}  $ is nested decreasing
according to $S(l,\mathbf{i},0)\supseteq S(l,\mathbf{i},1)\supseteq\dots$ and
there is finite $K$ such that the sequence converges in a finite number of
steps to a set $S(l,\mathbf{i})$ with non-empty interior%
\[
S(l,\mathbf{i}):=S(l,\mathbf{i},K)=S(l,\mathbf{i},K+1)=\dots.
\]


(ii) Let%
\[
S(l):=\{S(l,\mathbf{i}):\mathbf{i}\in\Omega_{l}\}.
\]
Then $\left\{  S(l)\right\}  _{l=0}^{\infty}$ is a nested increasing sequence
of collections of sets with non-empty interiors, according to
\[
S(0)\subset S(1)\subset\dots.
\]


(iii) Define $\Pi(\mathbf{j})$ as the union limit of the above sequence, that
is
\[
\Pi(\mathbf{j}):=%
%TCIMACRO{\tbigcup \limits_{n=1}^{\infty}}%
%BeginExpansion
{\textstyle\bigcup\limits_{n=1}^{\infty}}
%EndExpansion
S(n).
\]
Then $\Pi(\mathbf{j})$ provides a tiling of a region that includes
\[
f_{j_{1}}^{-1}\circ f_{j_{2}}^{-1}\circ\dots\circ f_{j_{l}}^{-1}\circ
f_{j_{l+1}}^{-1}(A_{j_{l+1}})
\]
for all sufficiently large $l$. In particular, if $\mathbf{j}$ includes the
symbol $1$ infinitely often, then $\Pi(\mathbf{j})$ provides a tiling of a
region that includes the blowup $\mathcal{A}\left(  \mathbf{j}\right)  $.
\end{theorem}

\begin{proof}
We begin by noting the following. The sets $A_{j_{k}\dots j_{l+2}j_{l+1}%
i_{1}\dots i_{l+1}}$, as in the preamble to the theorem, are non-empty by
Lemmas 2 and 3. Each of these sets has non-empty interior, by the following
argument. $A$ has non-empty interior by assumption. Because $A$ is an
attractor of an IFS\ of invertible maps, it is the closure of its interior,
i.e. $A=\overline{(A^{\circ})}$. Since the composition of homeomorphisms
preserves closure and non-empty interior, by Equation \ref{coreydef} in
Section 2.3, $A_{j_{k}\dots j_{l+2}j_{l+1}i_{1}\dots i_{l+1}}$ is of the form
$E\backslash G$ where $E$ has non-empty interior and $G$ is closed, so
inductively $A_{j_{k}\dots j_{l+2}j_{l+1}i_{1}\dots i_{l+1}}$ has non-empty
interior by Lemma \ref{lemmafour}.

To prove the first part of (i), observe that by Lemma \ref{topshiftlemma},
\[
f_{j_{k+1}}(A_{j_{k}\dots j_{l+2}j_{l+1}i_{1}\dots i_{l+1}})\supseteq
A_{j_{k+1}j_{k}\dots j_{l+2}j_{l+1}i_{1}\dots i_{l+1}}%
\]
and apply $f_{j_{1}}^{-1}\dots f_{j_{k+1}}^{-1}$ to both sides.

For the remainder of (i), notice that if $f_{j_{k+1}}(A_{j_{k}\dots
j_{l+2}j_{l+1}i_{1}\dots i_{l+1}})\cap C=\emptyset,$ then for any
$f_{\omega_{1}\omega_{2}\dots\omega_{k+1}}(A)$ intersecting $f_{j_{k+1}%
}(A_{j_{k}\dots j_{l+2}j_{l+1}i_{1}\dots i_{l+1}})$ such that $\omega
_{1}\omega_{2}\dots\omega_{k+1}>j_{k+1}j_{k}\dots j_{l+2}j_{l+1}i_{1}\dots
i_{l+1}$, we must have $\omega_{1}=j_{k+1}$. That is, $f_{j_{k+1}}%
(A_{j_{k}\dots j_{l+2}j_{l+1}i_{1}\dots i_{l+1}})=A_{j_{k+1}j_{k}\dots
j_{l+2}j_{l+1}i_{1}\dots i_{l+1}},$ and by applying $f_{j_{1}}^{-1}f_{j_{2}%
}^{-1}\dots f_{j_{k+1}}^{-1}$ to both sides we get $S(l,\mathbf{i,}%
k)=S(l,\mathbf{i,}k+1).$ So it suffices to find a large enough $K$ such that
$A_{j_{k}\dots j_{l+2}j_{l+1}i_{1}\dots i_{l+1}}$ contains no points of $C$
for all $k\geq K$.

But $A_{j_{k}\dots j_{l+2}j_{l+1}i_{1}\dots i_{l+1}}$ is contained in
$f_{j_{k+1}j_{k}\dots j_{l+2}j_{l+1}}(A)$ and for any fixed $l$ and $k$
sufficiently larger than $l,$ the condition on $\mathbf{j\in}%
\overleftarrow{\Sigma}$, $\widehat{\lambda}\in(\lambda,1),$ $x_{0}\in A,$ and
$\varepsilon>0,$ ensures that $f_{j_{k+1}j_{k}\dots j_{l+2}j_{l+1}}(A)$ does
not meet $C.$ Indeed,%
\begin{equation}
d(f_{j_{k+1}j_{k}\dots j_{l+2}j_{l+1}l+1\dots j_{1}}(x_{0}),f_{j_{k+1}%
j_{k}\dots j_{l+2}j_{l+1}}(A))\leq\lambda^{k-l}|A|, \label{estimate1}%
\end{equation}
where $\left\vert A\right\vert =\max\left\{  d(x,y)|x,y\in A\right\}  $ with
$d$ the metric on $\mathbb{R}^{n},$ and%
\begin{equation}
d(f_{j_{k+1}j_{k}\dots j_{l+2}j_{l+1}\dots j_{1}}(x_{0}),C)\geq\left(
\widehat{\lambda}\right)  ^{k+1}\varepsilon. \label{estimate2}%
\end{equation}
So the shortest distance between $f_{j_{k+1}j_{k}\dots j_{l+2}j_{l+1}}(A)$ and
$C$ is greater than
\begin{equation}
\left(  \widehat{\lambda}\right)  ^{k+1}\varepsilon-\lambda^{k-l}|A|
\label{ineq}%
\end{equation}
which is strictly positive, for any fixed $l$ and any $k$ that is sufficienly
larger than $l.$

To prove (ii), namely that $S(l)\subset S(l+1),$ note that
\begin{align*}
S(l)  &  =\{S(l,\mathbf{i}):\mathbf{i}\in\Omega_{l}\}\\
&  =\{f_{j_{1}}^{-1}\dots f_{j_{k}}^{-1}(A_{j_{k}\dots j_{l+2}j_{l+1}%
i_{1}\dots i_{l+1}}):\text{for all }k\text{ as in (i)},\text{ }\mathbf{i}%
\in\Omega_{l}\},
\end{align*}
while
\begin{align*}
S(l+1)  &  =\{S(l+1,\mathbf{i}):\mathbf{i}\in\Omega_{l+1}\}\\
&  =\{f_{j_{1}}^{-1}\dots f_{j_{k}}^{-1}(A_{j_{k}\dots j_{l+3}j_{l+2}%
i_{1}\dots i_{l+2}}):\text{for all }k\text{ as in (i)},\mathbf{i}\in
\Omega_{l+1}\}.
\end{align*}
Fix $k$ sufficiently large, say $k\geq K,$ so that both expressions on the
right hold. Then, using Lemma \ref{smallerlemma}, the set of allowed indices
in the first expression, ones with tails of the form $j_{l+1}i_{1}\dots
i_{l+1},$ where $l+1$ is fixed, is a subset of the set of allowed indices in
the second expression, with tails of the form $i_{1}\dots i_{l+2}$. So
$S(l)\subset S(l+1)$. This completes the proof of (ii).

To prove (iii) we make the following observations. (a) Let $K_{l}$ be the
maximum value of
\[
K_{l}:=\min\{k>l:f_{j_{k+1}j_{k}\dots j_{l+2}j_{l+1}i_{1}\dots i_{l+1}}(A)\cap
C=\varnothing\text{ }\forall\text{ }\mathbf{i}\in\Omega_{l}\}.
\]
Then it follows from (i) that the $k>l+1$ in the definition of $\Omega_{l}$
can be replaced by $l+1\leq k\leq K_{l}.$ (b) It can be seen from Equation
\ref{ineq} that $\left\{  K_{l}\right\}  $ is a non-increasing. In fact, upon
expansion, for%
\begin{equation}
l>(\log_{\widehat{\lambda}/\lambda}(|A|/\varepsilon\lambda)-1)/(1+\log
_{\widehat{\lambda}/\lambda}(\lambda)) \label{estimate3}%
\end{equation}
we have $K_{l}=l+1$.

It follows that
\begin{align*}%
%TCIMACRO{\tbigcup }%
%BeginExpansion
{\textstyle\bigcup}
%EndExpansion
S(l)  &  =%
%TCIMACRO{\tbigcup \limits_{i\in\Omega_{l}}^{\infty}}%
%BeginExpansion
{\textstyle\bigcup\limits_{i\in\Omega_{l}}^{\infty}}
%EndExpansion
f_{j_{1}}^{-1}\dots f_{j_{K_{l}}}^{-1}(A_{j_{K_{l}}\dots j_{l+1}i_{1}\dots
i_{l+1}})\\
&  =f_{-\mathbf{j}|l}(f_{-\sigma^{l}(\mathbf{j)|}K_{l}-l}(A_{j_{K_{l}}\dots
j_{l+1}}))\text{ by (a)}\\
&  =f_{-\mathbf{j}|l}(f_{j_{l+1}}^{-1}(A_{j_{l+1}}))\text{ for all large
enough }l\text{, as in Equation \ref{estimate3}.}%
\end{align*}
The penultimate statement in the theorem follows on noting that $f_{1}%
^{-1}(A_{1})=A$ and that $\left\{  f_{-\mathbf{j}|l}(A)\right\}
_{l=1}^{\infty}$ is a nested increasing sequence. In particular, the sequence
contains $A,$ and in particular the fixed-point of $f_{1},$ for all large
enough $l,$ and is acted on by $f_{1}^{-1}$ infinitely many times, so it
contains $f_{1}^{-1}(A),$ $f_{1}^{-1}\circ f_{1}^{-1}(A),\dots.$
\end{proof}

\end{document}
